\section{Controller Design}\label{sec:control}

\subsection{Tracking objective and timing}
Let $q_d\in C^2([0,\infty);\mathbb R^n)$ be a bounded reference with bounded
derivatives and define
\begin{equation}
 e=q-q_d,\qquad \dot e=\dot q-\dot q_d,\qquad z=[e^\top,\dot e^\top]^\top .
 \label{eq:errors}
\end{equation}
The total settling time is prescribed by two design parameters, the
observation time $T_o$ and the delay $h$:
\begin{equation}
 T=T_o+2h .
 \label{eq:prescribed}
\end{equation}
The objective is $e(t)=\dot e(t)=0$ for $t\ge T$ under exact compensation,
and an explicit ultimate bound in terms of the residual acceleration otherwise.

\subsection{Prescribed-time disturbance observer}\label{sec:ptdo}
Let $\chi=\dot q$ and write the mechanical subsystem as
\begin{equation}
 \dot\chi=u_o+\Delta_a,\qquad
 u_o=M_e^{-1}\bigl[\tau-C_e\dot q\bigr],\qquad
 \Delta_a=M_e^{-1}d .
 \label{eq:channel}
\end{equation}
The following observer is used:
\begin{equation}
 \dot z_1=z_2+u_o+\ell_1(t)\phi_1(\varepsilon_1/\sigma),\qquad
 \dot z_2=\ell_2(t)\phi_2(\varepsilon_1/\sigma),
 \label{eq:ptdo}
\end{equation}
with $\varepsilon_1=\chi-z_1$,
\begin{equation}
\begin{aligned}
 \phi_1(x)&=\operatorname{sig}^{1-\eta/2}(x)+\operatorname{sig}^{1+\eta/2}(x),\\
 \phi_2(x)&=\operatorname{sig}^{1-\eta}(x)+\operatorname{sig}^{1+\eta}(x),
\end{aligned}
 \label{eq:phi}
\end{equation}
$\operatorname{sig}^{\alpha}(x)=|x|^{\alpha}\operatorname{sgn}(x)$,
$0<\eta<1$, and the prescribed-time gains
\begin{equation}
 \ell_1(t)=\frac{m_1}{T_o-t},\qquad
 \ell_2(t)=\frac{m_2}{(T_o-t)^2},\qquad t<T_o .
 \label{eq:gains}
\end{equation}
The powers in \eqref{eq:phi} are dilation-consistent,
$1-\eta=2(1-\eta/2)-1$ and $1+\eta=2(1+\eta/2)-1$, so each branch of the error
system is homogeneous of negative degree with weights $(1,1\mp\eta/2)$; the
gains \eqref{eq:gains} serve as the time rescaling $s=-\ln(1-t/T_o)$.

\begin{assumption}[Observer interface]\label{as:interface}
There exist $T_o>0$ and an estimate $\hat\Delta_a$ such that, under the
continuous-time observer \eqref{eq:ptdo}, $\hat\Delta_a(t)=\Delta_a(t)$ for all
$t\ge T_o$, with a bounded error before $T_o$. The disturbance estimate fed to
the controller is $\hat d=M_e\hat\Delta_a$.
\end{assumption}

\begin{remark}\label{rem:frozen}
In the implementation the gains \eqref{eq:gains} are frozen at their value at
$T_o-\delta$. The frozen observer is a high-gain second-order differentiator
with residual $O(\|\dot\Delta_a\|\delta^2/m_2)$, which is part of the residual
$\Delta_r$ of Section~\ref{sec:stability}; it affects the size of the practical
bound, not the stability statement.
\end{remark}

\subsection{Dynamic compensation}
Using
\begin{equation}
 \tau=M_e(q)v+C_e(q,\dot q)\dot q-\hat d,
 \qquad v=\ddot q_d+\nu,
 \label{eq:computed_torque}
\end{equation}
and invoking Assumption~\ref{as:interface} for $t\ge T_o$, the error dynamics
of \eqref{eq:reduced} become the double integrator
\begin{equation}
 \dot z=Az+B\nu,\qquad
 A=\begin{bmatrix}0&I\\0&0\end{bmatrix},\quad
 B=\begin{bmatrix}0\\I\end{bmatrix},
 \label{eq:error_channel}
\end{equation}
which is controllable since $\mathrm{rank}[B,AB]=2n$.

\subsection{Smooth periodic delayed feedback}\label{sec:pdf}
Choose $K_0$ so that $A_c=A+BK_0$ is Hurwitz, let $h>0$, and let
$R_h:[0,2h]\to\mathbb R$ be smooth, nonnegative, zero on $[0,h]$ and positive
on a subset of positive measure in $(h,2h)$. Define the Gramian
\begin{equation}
 W_c=\int_h^{2h}e^{-A_cs}BR_h(s)B^\top e^{-A_c^\top s}\,ds
 \label{eq:wc}
\end{equation}
and, for the local time $\vartheta=t-T_o$, the $2h$-periodic gain
\begin{equation}
 K_c(\vartheta)=
 \begin{cases}
 0, & 0\le\vartheta<h,\\
 R_h(\vartheta)B^\top e^{-A_c^\top\vartheta}W_c^{-1}
 e^{A_c(h-\vartheta)}, & h\le\vartheta<2h .
 \end{cases}
 \label{eq:kc}
\end{equation}
The control input is
\begin{equation}
 \nu(t)=K_0z(t)-K_c(t-T_o)z(t-h),\qquad t\ge T_o,
 \label{eq:nu}
\end{equation}
so that the complete law is
\begin{equation}
 \tau=M_e\bigl[\ddot q_d+K_0z-K_c(t-T_o)z(t-h)\bigr]+C_e\dot q-\hat d .
 \label{eq:law}
\end{equation}
During the first half-period $K_c=0$ and the delayed state is only recorded;
the delayed term acts in the second half-period, where its product with the
phase-one trajectory cancels the period map.

\begin{remark}[Gain selection]\label{rem:gains}
A uniform baseline gain $K_0=[-K_p,\,-K_d]$ with $K_p=k_pI_n$,
$K_d=k_dI_n$ is used. Then $A_c$ is block diagonal with identical $2\times2$
blocks and $W_c$ inherits the same structure, with a condition number of
$7.4\times10^3$ for the example below. Normalising $K_p$ by the joint inertia
equalises the torque per radian of error but spreads the modal time scales over
three decades; in that case $\mathrm{cond}(W_c)$ exceeds $10^8$ and the
one-period nullification is destroyed by the inversion in \eqref{eq:kc}. The
uniform choice is therefore both numerically sound and, for this actuator
class, physically necessary.
\end{remark}
